% !TeX program = lualatex
% !TeX root = predavanje11.tex
% ============================================================
% Zajednički preamble za sva predavanja RWA kolegija
% Uključiti s: % ============================================================
% Zajednički preamble za sva predavanja RWA kolegija
% Uključiti s: % ============================================================
% Zajednički preamble za sva predavanja RWA kolegija
% Uključiti s: \input{../common/preamble.tex}
% ============================================================

\documentclass[aspectratio=149,11pt]{beamer}
\usetheme{CambridgeUS}

\usepackage{fontspec}              % radi samo sa XeLaTeX/LuaLaTeX
\usepackage[croatian]{babel}
\usepackage{csquotes}              % \enquote{} → hrvatski navodnici „..."
\usepackage{booktabs}
\usepackage{array}
\usepackage{listings}
\usepackage{xcolor}
\usepackage{colortbl}
\usepackage{hyperref}
\usepackage{tikz}
\usetikzlibrary{positioning,arrows.meta}

% ---------- Tipografija: IBM Plex ----------
\setsansfont{IBMPlexSans}[
  Extension      = .otf,
  UprightFont    = *-Regular,
  ItalicFont     = *-Italic,
  BoldFont       = *-SemiBold,
  BoldItalicFont = *-SemiBoldItalic
]
\setmainfont{IBMPlexSans}[
  Extension      = .otf,
  UprightFont    = *-Regular,
  ItalicFont     = *-Italic,
  BoldFont       = *-SemiBold,
  BoldItalicFont = *-SemiBoldItalic
]
\setmonofont{IBMPlexMono}[
  Extension      = .otf,
  UprightFont    = *-Regular,
  ItalicFont     = *-Italic,
  BoldFont       = *-SemiBold,
  BoldItalicFont = *-SemiBoldItalic,
  Scale          = 0.92
]

% ---------- Paleta ----------
\definecolor{ink}{HTML}{0C4685}      % primarna tamnoplava
\definecolor{ink2}{HTML}{6E7BAA}     % sekundarna plavo-siva
\definecolor{tint}{HTML}{EAEEF6}     % svjetli plavi tint
\definecolor{textclr}{HTML}{1A222E}  % glavni tekst
\definecolor{mutedclr}{HTML}{6B7484} % podnožje, sekundarni tekst
\definecolor{bgclr}{HTML}{F6F7F9}    % pozadina slajda
\definecolor{lineclr}{HTML}{D9DEE7}  % tanke linije tablica
\definecolor{divrule}{HTML}{9DB4D6}  % akcent na tamnim sekcijskim slajdovima
\definecolor{divtext}{HTML}{E6ECF6}  % tekst na tamnim sekcijskim slajdovima

\setbeamercolor{normal text}{fg=textclr,bg=bgclr}
\setbeamercolor{structure}{fg=ink}
\setbeamercolor{alerted text}{fg=ink}

\setbeamercolor{title}{fg=ink,bg=tint}
\setbeamercolor{subtitle}{fg=mutedclr}
\setbeamercolor{author}{fg=textclr}
\setbeamercolor{date}{fg=mutedclr}
\setbeamercolor{frametitle}{fg=ink,bg=tint}

\setbeamercolor{section in head/foot}{fg=white,bg=ink}
\setbeamercolor{subsection in head/foot}{fg=ink,bg=tint}
\setbeamercolor{author in head/foot}{fg=white,bg=ink}
\setbeamercolor{title in head/foot}{fg=textclr,bg=tint}
\setbeamercolor{date in head/foot}{fg=mutedclr,bg=tint}

\setbeamercolor{block title}{fg=white,bg=ink}
\setbeamercolor{block body}{fg=textclr,bg=tint}
\setbeamercolor{block title example}{fg=white,bg=ink2}
\setbeamercolor{block body example}{fg=textclr,bg=tint}
\setbeamercolor{itemize item}{fg=ink2}
\setbeamercolor{itemize subitem}{fg=ink2}
\setbeamercolor{enumerate item}{fg=ink2}

\setbeamertemplate{itemize items}[square]
\setbeamertemplate{navigation symbols}{}
\arrayrulecolor{lineclr}

% ---------- Kolegij, autori i logo ----------
\newcommand{\kolegij}{Programiranje za web}
\author{doc.\,dr.\,sc. Tomislav Petković \and mag.\,ing. Josip Torić}
\newcommand{\znak}[2]{\raisebox{-0.5\height}{\includegraphics[height=#1]{#2}}}
\titlegraphic{\znak{2.1cm}{../common/SIT_logo_tekst1.pdf}\hspace{1.8em}\znak{1.8cm}{../common/UNIZD_logo.pdf}}

% ---------- Zaglavlje ----------
% Jedna puna traka: lijevo sekcija, desno podsekcija.
\setbeamertemplate{headline}{%
  \leavevmode%
  \hbox{%
    \begin{beamercolorbox}[wd=\paperwidth,ht=2.25ex,dp=1ex,leftskip=.3cm,rightskip=.3cm]{section in head/foot}%
      \usebeamerfont{section in head/foot}\insertsectionhead\hfill\insertsubsectionhead%
    \end{beamercolorbox}%
  }%
}

% ---------- Podnožje ----------
% Lijevo: naziv trenutne prezentacije (kratki oblik naslova iz \title[...]{...})
% Sredina: naziv kolegija.  Desno: datum i redni broj slajda.
\setbeamertemplate{footline}{%
  \leavevmode%
  \hbox{%
    \begin{beamercolorbox}[wd=.34\paperwidth,ht=2.25ex,dp=1ex,center]{author in head/foot}%
      \usebeamerfont{author in head/foot}\insertshorttitle%
    \end{beamercolorbox}%
    \begin{beamercolorbox}[wd=.33\paperwidth,ht=2.25ex,dp=1ex,center]{title in head/foot}%
      \usebeamerfont{title in head/foot}\kolegij%
    \end{beamercolorbox}%
    \begin{beamercolorbox}[wd=.33\paperwidth,ht=2.25ex,dp=1ex,center]{date in head/foot}%
      \usebeamerfont{date in head/foot}\insertshortdate{}\hspace*{2em}%
      \insertframenumber{}\,/\,\inserttotalframenumber\hspace*{1ex}%
    \end{beamercolorbox}%
  }%
  \vskip0pt%
}

% ---------- Sekcijski slajd ----------
% \sekcija{Naslov sekcije}{jedna linija koja nabraja teme sekcije}
\newcommand{\sekcija}[2]{%
  \section{#1}%
  {%
    \setbeamercolor{background canvas}{bg=ink}%
    \setbeamertemplate{headline}{}%
    \begin{frame}
      \vfill
      \hspace*{0.03\textwidth}{\color{divrule}\rule{2.6cm}{1.1pt}}\par
      \vskip 0.7em
      \hspace*{0.03\textwidth}{\Large\bfseries\color{white}#1}\par
      \vskip 0.6em
      \hspace*{0.03\textwidth}\begin{minipage}{0.9\textwidth}\color{divtext}\small #2\end{minipage}\par
      \vfill
    \end{frame}%
  }%
}

% ---------- Listings ----------
\lstdefinelanguage{bash}{
  morekeywords={sudo,apt,brew,pip,pipx,uv,winget,python,git,flask,source,deactivate,code,gunicorn,docker,compose,heroku,render,railway,node,npm,pnpm,yarn,uvicorn,alembic,pytest,ruff,black,precommit},
  sensitive=false,
  morecomment=[l]{\#}
}
\lstdefinelanguage{python}{
  morekeywords={from,import,def,class,return,if,elif,else,for,in,while,try,except,raise,with,as,not,and,or,is,lambda,yield,True,False,None,await,async},
  sensitive=true,
  morecomment=[l]{\#},
  morestring=[b]',
  morestring=[b]"
}
\lstdefinelanguage{sql}{
  morekeywords={SELECT,FROM,WHERE,INSERT,INTO,VALUES,UPDATE,SET,DELETE,CREATE,TABLE,PRIMARY,KEY,FOREIGN,REFERENCES,UNIQUE,NOT,NULL,INDEX,ON,ALTER,DROP,CONSTRAINT},
  sensitive=false,
  morecomment=[l]{--}
}
\lstdefinelanguage{json}{
  morecomment=[l]{//},
  morestring=[b]",
  showstringspaces=false
}
\lstdefinelanguage{typescript}{
  morekeywords={import,export,from,const,let,var,function,return,if,else,for,of,in,while,switch,case,break,async,await,try,catch,finally,throw,new,class,this,extends,implements,interface,type,enum,readonly,public,private,protected,static,get,set,true,false,null,undefined,void,string,number,boolean,any,unknown,never,typeof,keyof,as},
  sensitive=true,
  morecomment=[l]{//},
  morecomment=[s]{/*}{*/},
  morestring=[b]',
  morestring=[b]",
  morestring=[b]`
}
\lstalias{javascript}{typescript}
\lstset{
  backgroundcolor=\color{tint},
  basicstyle=\ttfamily\footnotesize,
  keywordstyle=\color{ink}\bfseries,
  commentstyle=\color{mutedclr},
  frame=single,
  rulecolor=\color{lineclr},
  frameround=tttt,
  breaklines=true,
  tabsize=2,
  showstringspaces=false,
  upquote=true,
  columns=fullflexible,
  escapeinside={<@}{@>}
}
% ---------- Dijakritici u isječcima koda ----------
% listings čita bajt po bajt, pa višebajtni znak premjesti: "početni" ispadne
% "čpoetni". literate to ne popravlja. Rješenje je izlaz u LaTeX: sve između
% <@ i @> ne prolazi kroz listings, nego se slaže kao običan tekst, a unutar
% okvira je font ionako ttfamily, pa se ništa ne razlikuje.
%
%   async def get(...):
%       club = await repo.get_by_id(db, club_id)   <@\kom{dohvat kluba}@>
%       raise AppError("not_found", <@"Klub nije pronađen"@>, 404)
%
% Graničnici <@ i @> odabrani su jer se ne pojavljuju ni u jednom jeziku
% koji se na predavanjima citira; (* bi se sudaralo s def f(*args).
\newcommand{\kom}[1]{{\color{mutedclr}\# #1}}
% Isto za TypeScript, Vue i CSS, kojima komentar počinje s //:
%   const lifters = useLiftersStore()   <@\komts{dijeljeni popis}@>
\newcommand{\komts}[1]{{\color{mutedclr}// #1}}

% ---------- Činjenica naspram konvencije ----------
% Sve što nije propisano standardom ili jezikom, nego je dogovor struke,
% ide u ovaj okvir, da student zna koje je vrste tvrdnja pred njim.
% \begin{konvencija} ... \end{konvencija}
% \begin{konvencija}[Konvencija: nazivi putanja] ... \end{konvencija}
\newenvironment{konvencija}[1][Konvencija]
  {\begin{exampleblock}{#1}}
  {\end{exampleblock}}

% ---------- Helper macros ----------
% Stupac tablice zadane širine, poravnat ulijevo (bez rastezanja i dijeljenja riječi).
\newcolumntype{L}[1]{>{\raggedright\arraybackslash}p{#1}}
\newcommand{\codefile}[1]{\texttt{#1}}
\newcommand{\kbd}[1]{\texttt{#1}}

% ---------- Code frame environment ----------
% Automatski dodaje fragile i shrink=5, koristi se umjesto \begin{frame}
% Primjer: \begin{codeframe}{Naslov} ... \end{codeframe}
% Primjer s opcijama: \begin{codeframe}[shrink=15]{Naslov} ... \end{codeframe}
\newenvironment{codeframe}[2][]
  {\begin{frame}[fragile,shrink=5,#1]{#2}}
  {\end{frame}}

% ============================================================

\documentclass[aspectratio=149,11pt]{beamer}
\usetheme{CambridgeUS}

\usepackage{fontspec}              % radi samo sa XeLaTeX/LuaLaTeX
\usepackage[croatian]{babel}
\usepackage{csquotes}              % \enquote{} → hrvatski navodnici „..."
\usepackage{booktabs}
\usepackage{array}
\usepackage{listings}
\usepackage{xcolor}
\usepackage{colortbl}
\usepackage{hyperref}
\usepackage{tikz}
\usetikzlibrary{positioning,arrows.meta}

% ---------- Tipografija: IBM Plex ----------
\setsansfont{IBMPlexSans}[
  Extension      = .otf,
  UprightFont    = *-Regular,
  ItalicFont     = *-Italic,
  BoldFont       = *-SemiBold,
  BoldItalicFont = *-SemiBoldItalic
]
\setmainfont{IBMPlexSans}[
  Extension      = .otf,
  UprightFont    = *-Regular,
  ItalicFont     = *-Italic,
  BoldFont       = *-SemiBold,
  BoldItalicFont = *-SemiBoldItalic
]
\setmonofont{IBMPlexMono}[
  Extension      = .otf,
  UprightFont    = *-Regular,
  ItalicFont     = *-Italic,
  BoldFont       = *-SemiBold,
  BoldItalicFont = *-SemiBoldItalic,
  Scale          = 0.92
]

% ---------- Paleta ----------
\definecolor{ink}{HTML}{0C4685}      % primarna tamnoplava
\definecolor{ink2}{HTML}{6E7BAA}     % sekundarna plavo-siva
\definecolor{tint}{HTML}{EAEEF6}     % svjetli plavi tint
\definecolor{textclr}{HTML}{1A222E}  % glavni tekst
\definecolor{mutedclr}{HTML}{6B7484} % podnožje, sekundarni tekst
\definecolor{bgclr}{HTML}{F6F7F9}    % pozadina slajda
\definecolor{lineclr}{HTML}{D9DEE7}  % tanke linije tablica
\definecolor{divrule}{HTML}{9DB4D6}  % akcent na tamnim sekcijskim slajdovima
\definecolor{divtext}{HTML}{E6ECF6}  % tekst na tamnim sekcijskim slajdovima

\setbeamercolor{normal text}{fg=textclr,bg=bgclr}
\setbeamercolor{structure}{fg=ink}
\setbeamercolor{alerted text}{fg=ink}

\setbeamercolor{title}{fg=ink,bg=tint}
\setbeamercolor{subtitle}{fg=mutedclr}
\setbeamercolor{author}{fg=textclr}
\setbeamercolor{date}{fg=mutedclr}
\setbeamercolor{frametitle}{fg=ink,bg=tint}

\setbeamercolor{section in head/foot}{fg=white,bg=ink}
\setbeamercolor{subsection in head/foot}{fg=ink,bg=tint}
\setbeamercolor{author in head/foot}{fg=white,bg=ink}
\setbeamercolor{title in head/foot}{fg=textclr,bg=tint}
\setbeamercolor{date in head/foot}{fg=mutedclr,bg=tint}

\setbeamercolor{block title}{fg=white,bg=ink}
\setbeamercolor{block body}{fg=textclr,bg=tint}
\setbeamercolor{block title example}{fg=white,bg=ink2}
\setbeamercolor{block body example}{fg=textclr,bg=tint}
\setbeamercolor{itemize item}{fg=ink2}
\setbeamercolor{itemize subitem}{fg=ink2}
\setbeamercolor{enumerate item}{fg=ink2}

\setbeamertemplate{itemize items}[square]
\setbeamertemplate{navigation symbols}{}
\arrayrulecolor{lineclr}

% ---------- Kolegij, autori i logo ----------
\newcommand{\kolegij}{Programiranje za web}
\author{doc.\,dr.\,sc. Tomislav Petković \and mag.\,ing. Josip Torić}
\newcommand{\znak}[2]{\raisebox{-0.5\height}{\includegraphics[height=#1]{#2}}}
\titlegraphic{\znak{2.1cm}{../common/SIT_logo_tekst1.pdf}\hspace{1.8em}\znak{1.8cm}{../common/UNIZD_logo.pdf}}

% ---------- Zaglavlje ----------
% Jedna puna traka: lijevo sekcija, desno podsekcija.
\setbeamertemplate{headline}{%
  \leavevmode%
  \hbox{%
    \begin{beamercolorbox}[wd=\paperwidth,ht=2.25ex,dp=1ex,leftskip=.3cm,rightskip=.3cm]{section in head/foot}%
      \usebeamerfont{section in head/foot}\insertsectionhead\hfill\insertsubsectionhead%
    \end{beamercolorbox}%
  }%
}

% ---------- Podnožje ----------
% Lijevo: naziv trenutne prezentacije (kratki oblik naslova iz \title[...]{...})
% Sredina: naziv kolegija.  Desno: datum i redni broj slajda.
\setbeamertemplate{footline}{%
  \leavevmode%
  \hbox{%
    \begin{beamercolorbox}[wd=.34\paperwidth,ht=2.25ex,dp=1ex,center]{author in head/foot}%
      \usebeamerfont{author in head/foot}\insertshorttitle%
    \end{beamercolorbox}%
    \begin{beamercolorbox}[wd=.33\paperwidth,ht=2.25ex,dp=1ex,center]{title in head/foot}%
      \usebeamerfont{title in head/foot}\kolegij%
    \end{beamercolorbox}%
    \begin{beamercolorbox}[wd=.33\paperwidth,ht=2.25ex,dp=1ex,center]{date in head/foot}%
      \usebeamerfont{date in head/foot}\insertshortdate{}\hspace*{2em}%
      \insertframenumber{}\,/\,\inserttotalframenumber\hspace*{1ex}%
    \end{beamercolorbox}%
  }%
  \vskip0pt%
}

% ---------- Sekcijski slajd ----------
% \sekcija{Naslov sekcije}{jedna linija koja nabraja teme sekcije}
\newcommand{\sekcija}[2]{%
  \section{#1}%
  {%
    \setbeamercolor{background canvas}{bg=ink}%
    \setbeamertemplate{headline}{}%
    \begin{frame}
      \vfill
      \hspace*{0.03\textwidth}{\color{divrule}\rule{2.6cm}{1.1pt}}\par
      \vskip 0.7em
      \hspace*{0.03\textwidth}{\Large\bfseries\color{white}#1}\par
      \vskip 0.6em
      \hspace*{0.03\textwidth}\begin{minipage}{0.9\textwidth}\color{divtext}\small #2\end{minipage}\par
      \vfill
    \end{frame}%
  }%
}

% ---------- Listings ----------
\lstdefinelanguage{bash}{
  morekeywords={sudo,apt,brew,pip,pipx,uv,winget,python,git,flask,source,deactivate,code,gunicorn,docker,compose,heroku,render,railway,node,npm,pnpm,yarn,uvicorn,alembic,pytest,ruff,black,precommit},
  sensitive=false,
  morecomment=[l]{\#}
}
\lstdefinelanguage{python}{
  morekeywords={from,import,def,class,return,if,elif,else,for,in,while,try,except,raise,with,as,not,and,or,is,lambda,yield,True,False,None,await,async},
  sensitive=true,
  morecomment=[l]{\#},
  morestring=[b]',
  morestring=[b]"
}
\lstdefinelanguage{sql}{
  morekeywords={SELECT,FROM,WHERE,INSERT,INTO,VALUES,UPDATE,SET,DELETE,CREATE,TABLE,PRIMARY,KEY,FOREIGN,REFERENCES,UNIQUE,NOT,NULL,INDEX,ON,ALTER,DROP,CONSTRAINT},
  sensitive=false,
  morecomment=[l]{--}
}
\lstdefinelanguage{json}{
  morecomment=[l]{//},
  morestring=[b]",
  showstringspaces=false
}
\lstdefinelanguage{typescript}{
  morekeywords={import,export,from,const,let,var,function,return,if,else,for,of,in,while,switch,case,break,async,await,try,catch,finally,throw,new,class,this,extends,implements,interface,type,enum,readonly,public,private,protected,static,get,set,true,false,null,undefined,void,string,number,boolean,any,unknown,never,typeof,keyof,as},
  sensitive=true,
  morecomment=[l]{//},
  morecomment=[s]{/*}{*/},
  morestring=[b]',
  morestring=[b]",
  morestring=[b]`
}
\lstalias{javascript}{typescript}
\lstset{
  backgroundcolor=\color{tint},
  basicstyle=\ttfamily\footnotesize,
  keywordstyle=\color{ink}\bfseries,
  commentstyle=\color{mutedclr},
  frame=single,
  rulecolor=\color{lineclr},
  frameround=tttt,
  breaklines=true,
  tabsize=2,
  showstringspaces=false,
  upquote=true,
  columns=fullflexible,
  escapeinside={<@}{@>}
}
% ---------- Dijakritici u isječcima koda ----------
% listings čita bajt po bajt, pa višebajtni znak premjesti: "početni" ispadne
% "čpoetni". literate to ne popravlja. Rješenje je izlaz u LaTeX: sve između
% <@ i @> ne prolazi kroz listings, nego se slaže kao običan tekst, a unutar
% okvira je font ionako ttfamily, pa se ništa ne razlikuje.
%
%   async def get(...):
%       club = await repo.get_by_id(db, club_id)   <@\kom{dohvat kluba}@>
%       raise AppError("not_found", <@"Klub nije pronađen"@>, 404)
%
% Graničnici <@ i @> odabrani su jer se ne pojavljuju ni u jednom jeziku
% koji se na predavanjima citira; (* bi se sudaralo s def f(*args).
\newcommand{\kom}[1]{{\color{mutedclr}\# #1}}
% Isto za TypeScript, Vue i CSS, kojima komentar počinje s //:
%   const lifters = useLiftersStore()   <@\komts{dijeljeni popis}@>
\newcommand{\komts}[1]{{\color{mutedclr}// #1}}

% ---------- Činjenica naspram konvencije ----------
% Sve što nije propisano standardom ili jezikom, nego je dogovor struke,
% ide u ovaj okvir, da student zna koje je vrste tvrdnja pred njim.
% \begin{konvencija} ... \end{konvencija}
% \begin{konvencija}[Konvencija: nazivi putanja] ... \end{konvencija}
\newenvironment{konvencija}[1][Konvencija]
  {\begin{exampleblock}{#1}}
  {\end{exampleblock}}

% ---------- Helper macros ----------
% Stupac tablice zadane širine, poravnat ulijevo (bez rastezanja i dijeljenja riječi).
\newcolumntype{L}[1]{>{\raggedright\arraybackslash}p{#1}}
\newcommand{\codefile}[1]{\texttt{#1}}
\newcommand{\kbd}[1]{\texttt{#1}}

% ---------- Code frame environment ----------
% Automatski dodaje fragile i shrink=5, koristi se umjesto \begin{frame}
% Primjer: \begin{codeframe}{Naslov} ... \end{codeframe}
% Primjer s opcijama: \begin{codeframe}[shrink=15]{Naslov} ... \end{codeframe}
\newenvironment{codeframe}[2][]
  {\begin{frame}[fragile,shrink=5,#1]{#2}}
  {\end{frame}}

% ============================================================

\documentclass[aspectratio=149,11pt]{beamer}
\usetheme{CambridgeUS}

\usepackage{fontspec}              % radi samo sa XeLaTeX/LuaLaTeX
\usepackage[croatian]{babel}
\usepackage{csquotes}              % \enquote{} → hrvatski navodnici „..."
\usepackage{booktabs}
\usepackage{array}
\usepackage{listings}
\usepackage{xcolor}
\usepackage{colortbl}
\usepackage{hyperref}
\usepackage{tikz}
\usetikzlibrary{positioning,arrows.meta}

% ---------- Tipografija: IBM Plex ----------
\setsansfont{IBMPlexSans}[
  Extension      = .otf,
  UprightFont    = *-Regular,
  ItalicFont     = *-Italic,
  BoldFont       = *-SemiBold,
  BoldItalicFont = *-SemiBoldItalic
]
\setmainfont{IBMPlexSans}[
  Extension      = .otf,
  UprightFont    = *-Regular,
  ItalicFont     = *-Italic,
  BoldFont       = *-SemiBold,
  BoldItalicFont = *-SemiBoldItalic
]
\setmonofont{IBMPlexMono}[
  Extension      = .otf,
  UprightFont    = *-Regular,
  ItalicFont     = *-Italic,
  BoldFont       = *-SemiBold,
  BoldItalicFont = *-SemiBoldItalic,
  Scale          = 0.92
]

% ---------- Paleta ----------
\definecolor{ink}{HTML}{0C4685}      % primarna tamnoplava
\definecolor{ink2}{HTML}{6E7BAA}     % sekundarna plavo-siva
\definecolor{tint}{HTML}{EAEEF6}     % svjetli plavi tint
\definecolor{textclr}{HTML}{1A222E}  % glavni tekst
\definecolor{mutedclr}{HTML}{6B7484} % podnožje, sekundarni tekst
\definecolor{bgclr}{HTML}{F6F7F9}    % pozadina slajda
\definecolor{lineclr}{HTML}{D9DEE7}  % tanke linije tablica
\definecolor{divrule}{HTML}{9DB4D6}  % akcent na tamnim sekcijskim slajdovima
\definecolor{divtext}{HTML}{E6ECF6}  % tekst na tamnim sekcijskim slajdovima

\setbeamercolor{normal text}{fg=textclr,bg=bgclr}
\setbeamercolor{structure}{fg=ink}
\setbeamercolor{alerted text}{fg=ink}

\setbeamercolor{title}{fg=ink,bg=tint}
\setbeamercolor{subtitle}{fg=mutedclr}
\setbeamercolor{author}{fg=textclr}
\setbeamercolor{date}{fg=mutedclr}
\setbeamercolor{frametitle}{fg=ink,bg=tint}

\setbeamercolor{section in head/foot}{fg=white,bg=ink}
\setbeamercolor{subsection in head/foot}{fg=ink,bg=tint}
\setbeamercolor{author in head/foot}{fg=white,bg=ink}
\setbeamercolor{title in head/foot}{fg=textclr,bg=tint}
\setbeamercolor{date in head/foot}{fg=mutedclr,bg=tint}

\setbeamercolor{block title}{fg=white,bg=ink}
\setbeamercolor{block body}{fg=textclr,bg=tint}
\setbeamercolor{block title example}{fg=white,bg=ink2}
\setbeamercolor{block body example}{fg=textclr,bg=tint}
\setbeamercolor{itemize item}{fg=ink2}
\setbeamercolor{itemize subitem}{fg=ink2}
\setbeamercolor{enumerate item}{fg=ink2}

\setbeamertemplate{itemize items}[square]
\setbeamertemplate{navigation symbols}{}
\arrayrulecolor{lineclr}

% ---------- Kolegij, autori i logo ----------
\newcommand{\kolegij}{Programiranje za web}
\author{doc.\,dr.\,sc. Tomislav Petković \and mag.\,ing. Josip Torić}
\newcommand{\znak}[2]{\raisebox{-0.5\height}{\includegraphics[height=#1]{#2}}}
\titlegraphic{\znak{2.1cm}{../common/SIT_logo_tekst1.pdf}\hspace{1.8em}\znak{1.8cm}{../common/UNIZD_logo.pdf}}

% ---------- Zaglavlje ----------
% Jedna puna traka: lijevo sekcija, desno podsekcija.
\setbeamertemplate{headline}{%
  \leavevmode%
  \hbox{%
    \begin{beamercolorbox}[wd=\paperwidth,ht=2.25ex,dp=1ex,leftskip=.3cm,rightskip=.3cm]{section in head/foot}%
      \usebeamerfont{section in head/foot}\insertsectionhead\hfill\insertsubsectionhead%
    \end{beamercolorbox}%
  }%
}

% ---------- Podnožje ----------
% Lijevo: naziv trenutne prezentacije (kratki oblik naslova iz \title[...]{...})
% Sredina: naziv kolegija.  Desno: datum i redni broj slajda.
\setbeamertemplate{footline}{%
  \leavevmode%
  \hbox{%
    \begin{beamercolorbox}[wd=.34\paperwidth,ht=2.25ex,dp=1ex,center]{author in head/foot}%
      \usebeamerfont{author in head/foot}\insertshorttitle%
    \end{beamercolorbox}%
    \begin{beamercolorbox}[wd=.33\paperwidth,ht=2.25ex,dp=1ex,center]{title in head/foot}%
      \usebeamerfont{title in head/foot}\kolegij%
    \end{beamercolorbox}%
    \begin{beamercolorbox}[wd=.33\paperwidth,ht=2.25ex,dp=1ex,center]{date in head/foot}%
      \usebeamerfont{date in head/foot}\insertshortdate{}\hspace*{2em}%
      \insertframenumber{}\,/\,\inserttotalframenumber\hspace*{1ex}%
    \end{beamercolorbox}%
  }%
  \vskip0pt%
}

% ---------- Sekcijski slajd ----------
% \sekcija{Naslov sekcije}{jedna linija koja nabraja teme sekcije}
\newcommand{\sekcija}[2]{%
  \section{#1}%
  {%
    \setbeamercolor{background canvas}{bg=ink}%
    \setbeamertemplate{headline}{}%
    \begin{frame}
      \vfill
      \hspace*{0.03\textwidth}{\color{divrule}\rule{2.6cm}{1.1pt}}\par
      \vskip 0.7em
      \hspace*{0.03\textwidth}{\Large\bfseries\color{white}#1}\par
      \vskip 0.6em
      \hspace*{0.03\textwidth}\begin{minipage}{0.9\textwidth}\color{divtext}\small #2\end{minipage}\par
      \vfill
    \end{frame}%
  }%
}

% ---------- Listings ----------
\lstdefinelanguage{bash}{
  morekeywords={sudo,apt,brew,pip,pipx,uv,winget,python,git,flask,source,deactivate,code,gunicorn,docker,compose,heroku,render,railway,node,npm,pnpm,yarn,uvicorn,alembic,pytest,ruff,black,precommit},
  sensitive=false,
  morecomment=[l]{\#}
}
\lstdefinelanguage{python}{
  morekeywords={from,import,def,class,return,if,elif,else,for,in,while,try,except,raise,with,as,not,and,or,is,lambda,yield,True,False,None,await,async},
  sensitive=true,
  morecomment=[l]{\#},
  morestring=[b]',
  morestring=[b]"
}
\lstdefinelanguage{sql}{
  morekeywords={SELECT,FROM,WHERE,INSERT,INTO,VALUES,UPDATE,SET,DELETE,CREATE,TABLE,PRIMARY,KEY,FOREIGN,REFERENCES,UNIQUE,NOT,NULL,INDEX,ON,ALTER,DROP,CONSTRAINT},
  sensitive=false,
  morecomment=[l]{--}
}
\lstdefinelanguage{json}{
  morecomment=[l]{//},
  morestring=[b]",
  showstringspaces=false
}
\lstdefinelanguage{typescript}{
  morekeywords={import,export,from,const,let,var,function,return,if,else,for,of,in,while,switch,case,break,async,await,try,catch,finally,throw,new,class,this,extends,implements,interface,type,enum,readonly,public,private,protected,static,get,set,true,false,null,undefined,void,string,number,boolean,any,unknown,never,typeof,keyof,as},
  sensitive=true,
  morecomment=[l]{//},
  morecomment=[s]{/*}{*/},
  morestring=[b]',
  morestring=[b]",
  morestring=[b]`
}
\lstalias{javascript}{typescript}
\lstset{
  backgroundcolor=\color{tint},
  basicstyle=\ttfamily\footnotesize,
  keywordstyle=\color{ink}\bfseries,
  commentstyle=\color{mutedclr},
  frame=single,
  rulecolor=\color{lineclr},
  frameround=tttt,
  breaklines=true,
  tabsize=2,
  showstringspaces=false,
  upquote=true,
  columns=fullflexible,
  escapeinside={<@}{@>}
}
% ---------- Dijakritici u isječcima koda ----------
% listings čita bajt po bajt, pa višebajtni znak premjesti: "početni" ispadne
% "čpoetni". literate to ne popravlja. Rješenje je izlaz u LaTeX: sve između
% <@ i @> ne prolazi kroz listings, nego se slaže kao običan tekst, a unutar
% okvira je font ionako ttfamily, pa se ništa ne razlikuje.
%
%   async def get(...):
%       club = await repo.get_by_id(db, club_id)   <@\kom{dohvat kluba}@>
%       raise AppError("not_found", <@"Klub nije pronađen"@>, 404)
%
% Graničnici <@ i @> odabrani su jer se ne pojavljuju ni u jednom jeziku
% koji se na predavanjima citira; (* bi se sudaralo s def f(*args).
\newcommand{\kom}[1]{{\color{mutedclr}\# #1}}
% Isto za TypeScript, Vue i CSS, kojima komentar počinje s //:
%   const lifters = useLiftersStore()   <@\komts{dijeljeni popis}@>
\newcommand{\komts}[1]{{\color{mutedclr}// #1}}

% ---------- Činjenica naspram konvencije ----------
% Sve što nije propisano standardom ili jezikom, nego je dogovor struke,
% ide u ovaj okvir, da student zna koje je vrste tvrdnja pred njim.
% \begin{konvencija} ... \end{konvencija}
% \begin{konvencija}[Konvencija: nazivi putanja] ... \end{konvencija}
\newenvironment{konvencija}[1][Konvencija]
  {\begin{exampleblock}{#1}}
  {\end{exampleblock}}

% ---------- Helper macros ----------
% Stupac tablice zadane širine, poravnat ulijevo (bez rastezanja i dijeljenja riječi).
\newcolumntype{L}[1]{>{\raggedright\arraybackslash}p{#1}}
\newcommand{\codefile}[1]{\texttt{#1}}
\newcommand{\kbd}[1]{\texttt{#1}}

% ---------- Code frame environment ----------
% Automatski dodaje fragile i shrink=5, koristi se umjesto \begin{frame}
% Primjer: \begin{codeframe}{Naslov} ... \end{codeframe}
% Primjer s opcijama: \begin{codeframe}[shrink=15]{Naslov} ... \end{codeframe}
\newenvironment{codeframe}[2][]
  {\begin{frame}[fragile,shrink=5,#1]{#2}}
  {\end{frame}}


\title[Predavanje 11: Testiranje]{Programiranje za web \\ Predavanje 11: Testiranje cjelovite aplikacije}
\date{07.01.2027.}

\begin{document}

\frame{\titlepage}

% ------------------------------------------------------------
\begin{frame}{Plan predavanja}
\begin{enumerate}
  \item Što test dokazuje
  \item Testovi poslužitelja
  \item Testovi klijentskog dijela
  \item Test kroz preglednik
\end{enumerate}

\vspace{0.6em}
Aplikacija radi kad se isproba rukom. Ovo predavanje uređuje kako se to zna i nakon sljedeće izmjene: testom na razini na kojoj se pogreška može pojaviti, uz provjeru da test i mjera pokrivenosti vide ono što tvrde.
\end{frame}

% ------------------------------------------------------------
\begin{frame}{Ishodi učenja}
Nakon ovog predavanja student može:

\begin{itemize}
  \item odabrati razinu testa prema pogrešci koju test treba uhvatiti
  \item provjeriti test namjernim kvarom, a mjeru pokrivenosti retkom za koji zna da se izvodi
  \item napisati jedinični test klijentskog dijela koji ne ovisi o vremenskoj zoni ni o danu pokretanja
  \item zamijeniti modul ili funkciju preglednika lažnom inačicom u jediničnom testu
  \item napisati test kroz preglednik koji podatke priprema sam i elemente traži po ulozi i oznaci
  \item uključiti automatsku provjeru pristupačnosti u test kroz preglednik
\end{itemize}
\end{frame}

% ------------------------------------------------------------
\begin{frame}{Domena primjera}
Klubovi prijavljuju natjecatelje na natjecanja, a organizator unosi natjecanja s dvama rokovima. Backend (poslužiteljski dio) provodi pravila, a frontend (klijentski dio) ih prikazuje: ne nudi radnju koju rok više ne dopušta, ispisuje rokove po zagrebačkom vremenu, a u obrascu nudi samo kategorije koje odgovaraju spolu natjecatelja.

\vspace{0.9em}
Backend ima testove koji šalju zahtjeve nad zasebnom bazom. Frontend nema nijedan: sve što se o njemu zna, zna se jer je netko kliknuo i vidio da radi.

\vspace{0.9em}
Neki podatci postoje na objema stranama. Popis težinskih kategorija stoji i u shemi backenda i u kodu frontenda.
\end{frame}

% ------------------------------------------------------------
\begin{frame}[fragile]{Izmjena koja ništa ne sruši (I)}
Savez dopusti nastup juniorkama u kategoriji do 43 kg. Na backendu se popis ženskih kategorija proširi:

\begin{lstlisting}[language=python]
# app/schemas/registration.py
FEMALE_CATEGORIES = ["43", "47", "52", "57", "63", "69", ...]
\end{lstlisting}

\vspace{0.1em}
Na frontendu popis ostaje kakav je bio:

\begin{lstlisting}[language=typescript]
// src/lib/categories.ts
export const CATEGORIES: Record<Gender, readonly string[]> = {
  M: ['59', '66', '74', '83', '93', '105', '120', '120+'],
  F: ['47', '52', '57', '63', '69', '76', '84', '84+'],
}
\end{lstlisting}

\vspace{0.1em}
Testovi backenda prolaze, a frontend se gradi bez pogreške. Klub juniorku od 42 kg ne može prijaviti, jer obrazac ne nudi kategoriju 43. Dva popisa istog podatka razišla su se, a nijedna provjera to nije prijavila.
\end{frame}

% ------------------------------------------------------------
\begin{frame}{Izmjena koja ništa ne sruši (II)}
Prvi je odgovor komentar \enquote{mijenjati zajedno} uz oba popisa. U ovom projektu komentar stoji samo u datoteci frontenda, a i kad bi stajao na obama mjestima, čita ga samo onaj tko datoteku otvori.

\vspace{0.9em}
Drugi je odgovor jedan izvor: frontend dohvaća kategorije s backenda. Drugi popis time nestaje, ali API dobiva rutu više, a obrazac čeka na dohvat prije prikaza.

\vspace{0.9em}
Treći je odgovor provjera koja oba popisa pročita i usporedi. Izmjena koja ih razdvoji tada ne prolazi tiho: provjera pada s imenom pravila koje je prekršeno. Takva provjera je \textbf{test}: zapis pravila koji se izvodi.
\end{frame}

% ============================================================
\sekcija{Što test dokazuje}{Doseg prolaznog testa, razine testa i cijena svake, oblik i ime testa.}

% ------------------------------------------------------------
\begin{frame}{Doseg prolaznog testa}
Ispit pokazuje da student zna riješiti zadatke koji su bili na ispitu. O ostatku gradiva ne govori ništa, a zadatke je odabrao nastavnik.

\vspace{0.9em}
S testom je jednako. Kad prođe, pokazuje da je za ulaze koje je test zadao dobiven očekivan ishod. Ne pokazuje da je kod ispravan, jer ulaze bira onaj tko test piše.

\vspace{0.9em}
Odatle mjerilo izbora: testira se ono što se ne smije pokvariti, a ne ono što je lako napisati. Ruta ili prikaz bez testa ostaju neprovjereni, koliko god ostalih testova prolazilo.
\end{frame}

% ------------------------------------------------------------
\begin{frame}{Razine testa}
\vspace{0.3em}
\begin{center}
{\small
\setlength{\tabcolsep}{6pt}
\renewcommand{\arraystretch}{1.35}
\begin{tabular}{@{}L{2.6cm}L{4.3cm}L{4.4cm}@{}}
  \toprule
  \textbf{Razina} & \textbf{Obuhvaća} & \textbf{Cijena} \\
  \midrule
  jedinični & jednu funkciju ili modul, bez mreže i baze & ne vidi vezu među dijelovima \\
  nad API-jem & zahtjev kroz sve slojeve backenda do baze & traži bazu; ne vidi frontend \\
  kroz preglednik & put od klika do baze i natrag & traži pokrenuto sve; sporiji \\
  \bottomrule
\end{tabular}
}
\end{center}

\vspace{0.6em}
Niža razina pada brže i točnije: kad padne jedinični test, pogreška je u jednoj funkciji. Kad padne test kroz preglednik, pogreška je negdje između klika i baze.

\vspace{0.4em}
Testovi niže razine zato se pišu u velikom broju, a testovi kroz preglednik za nekoliko puteva kojima korisnik prolazi.
\end{frame}

% ------------------------------------------------------------
\begin{frame}{Oblik testa}
Svaki test ima tri dijela, uvijek istim redom:

\vspace{0.6em}
\begin{itemize}
  \itemsep0.5em
  \item \textbf{Priprema.} Stanje koje test traži: zapisi u bazi, prijavljeni račun, zadan trenutak.
  \item \textbf{Radnja.} Jedan poziv funkcije ili jedan zahtjev.
  \item \textbf{Tvrdnja.} Kakav je ishod morao biti: statusni kod, polje u odgovoru, tekst na zaslonu.
\end{itemize}

\vspace{0.8em}
Jedan test provjerava jedno pravilo, a ime mu je ime tog pravila: \texttt{test\_refresh\_rejected\_after\_deactivation}, a ne \texttt{test\_refresh\_works}. Kad test padne, ispis prvo kaže koje je pravilo prekršeno, a tvrdnja koja je pala kaže očekivanu i dobivenu vrijednost.
\end{frame}

% ============================================================
\sekcija{Testovi poslužitelja}{Mjera pokrivenosti, mjerni alat koji dio izvedenog koda ne vidi, test koji prolazi zbog druge provjere, granica mjere.}

% ------------------------------------------------------------
\begin{frame}[fragile]{Pokrivenost koda}
Nadzornik koji obilazi zgradu bilježi u koje je prostorije ušao. Popis kaže gdje nije bio, ali ne kaže je li u prostoriji u kojoj je bio išta pregledao.

\vspace{0.7em}
\textbf{Pokrivenost koda} (engl. \emph{code coverage}) bilježi isto za testove: koje je retke aplikacije izveo barem jedan test. Za pytest je mjeri dodatak \texttt{pytest-cov}:

\begin{lstlisting}[language=bash]
pytest --cov
\end{lstlisting}

\vspace{0.3em}
Ispis za svaku datoteku navodi broj naredbi (\texttt{Stmts}), broj neizvedenih (\texttt{Miss}), udio izvedenih (\texttt{Cover}) i brojeve neizvedenih redaka (\texttt{Missing}).
\end{frame}

% ------------------------------------------------------------
\begin{frame}[fragile]{Pokrivenost sloja services (I)}
Prvo mjerenje dalo je ovaj ispis, uz testove koji svi prolaze. Stupac \texttt{Missing} ovdje je izostavljen:

\begin{lstlisting}[keepspaces=true]
Name                                  Stmts   Miss  Cover
app/services/auth_service.py             31      9    71%
app/services/club_service.py             56     29    48%
app/services/competition_service.py      24     11    54%
app/services/lifter_service.py           35     17    51%
app/services/registration_service.py     78     44    44%
...
108 passed in 14.53s
\end{lstlisting}

\vspace{0.4em}
Prema mjeri, gotovo pola redaka sloja \codefile{services} nijedan test ne izvodi, a u tom sloju stoje poslovna pravila.
\end{frame}

% ------------------------------------------------------------
\begin{frame}{Pokrivenost sloja services (II)}
Prvi je odgovor dopisati testove za sloj \codefile{services}. Za retke koje mjera prijavljuje kao neizvedene testovi već postoje: test u kojem klub prijavljuje svojeg natjecatelja izvodi funkciju \texttt{create\_registration} i tvrdi da je odgovor \texttt{201}.

\vspace{0.9em}
Kad bi se ti retci doista preskakali, upis se ne bi dogodio i taj test ne bi prošao. Test prolazi, pa su se retci izveli.

\pause
\vspace{0.9em}
Pogreška je u mjernom alatu: dio izvedenog koda ne vidi. Prije nego se broju povjeruje, provjerava se redak za koji se zna da ga test izvodi. Ako ga mjera prijavljuje kao neizveden, kvar je u mjerenju.
\end{frame}

% ------------------------------------------------------------
\begin{frame}[fragile]{Postavka mjernog alata}
SQLAlchemy u asinkronom radu izvodi dio posla u \emph{greenletima}, lakim tokovima izvođenja koji se izmjenjuju unutar iste dretve operacijskog sustava.

\vspace{0.5em}
Bez posebne postavke alat \texttt{coverage.py} ne prati prelazak između greenleta, pa ne bilježi retke koji se izvedu nakon povratka iz baze, dakle iza \texttt{await}:

\begin{lstlisting}
# api/pyproject.toml
[tool.coverage.run]
source = ["app"]
concurrency = ["greenlet", "thread"]
omit = ["app/seed.py"]
\end{lstlisting}

\vspace{0.3em}
\texttt{concurrency} mjernom alatu kaže koje tokove izvođenja aplikacija koristi. \texttt{omit} izostavlja početne podatke, koji nisu dio aplikacije, pa ih testovi i ne izvode.
\end{frame}

% ------------------------------------------------------------
\begin{frame}[fragile]{Pokrivenost sloja services (III)}
Uz postavku iz prethodnog slajda isti testovi nad istom aplikacijom daju ovaj ispis:

\begin{lstlisting}[keepspaces=true]
Name                                  Stmts   Miss  Cover   Missing
app/services/auth_service.py             31      1    97%   48
app/services/club_service.py             56      3    95%   65-66, 87
app/services/competition_service.py      24      1    96%   40
app/services/lifter_service.py           35      2    94%   70-72
app/services/registration_service.py     78      2    97%   43, 69
\end{lstlisting}

\vspace{0.4em}
Datoteka \texttt{registration\_service.py} u prvom ispisu ima 44~\% izvedenih redaka, a u ovom 97~\%. Nijedan test nije promijenjen.

\vspace{0.4em}
Preostali retci sada su doista neizvedeni, a svaki od njih vrijedi pogledati.
\end{frame}

% ------------------------------------------------------------
\begin{frame}[fragile]{Test koji prolazi zbog druge provjere (I)}
Redak 43 u \texttt{registration\_service.py} odbija prijavu koja ne pripada natjecanju iz adrese. Za to pravilo test postoji i prolazi:

\begin{lstlisting}[language=python]
# tests/test_registrations.py
async def test_registration_under_wrong_competition_is_404(
    client, db, lifter, competition
):
    reg = await add_registration(db, lifter, competition, "83")
    headers = await auth_header(client, "testclub", "klub123")
    resp = await client.get(
        f"/competitions/9999/registrations/{reg.id}", headers=headers
    )
    assert resp.status_code == 404
\end{lstlisting}

\vspace{0.3em}
Mjera unatoč tome kaže da redak 43 nijedan test ne izvodi. Test nosi ime pravila do kojega ne dolazi.
\end{frame}

% ------------------------------------------------------------
\begin{frame}[fragile]{Test koji prolazi zbog druge provjere (II)}
\begin{lstlisting}[language=python]
# app/services/registration_service.py
async def get_registration(
    db: AsyncSession, comp_id: int, reg_id: int, current_user: User
) -> Registration:
    await get_competition_or_404(db, comp_id)
    reg = await _get_registration_or_404(db, comp_id, reg_id)
    _check_lifter_ownership(reg.lifter, current_user)
    return reg
\end{lstlisting}

\vspace{0.4em}
Natjecanje 9999 ne postoji, pa zahtjev odbija prva provjera, s istim statusom \texttt{404}. Do provjere pripada li prijava natjecanju iz adrese ne dolazi.

\vspace{0.4em}
Tvrdnja o polju \texttt{code} ne bi pomogla, jer obje provjere vraćaju \texttt{not\_found}. Test treba drugu pripremu: natjecanje iz adrese mora postojati. Tada zahtjev prolazi prvu provjeru, odbija ga tek provjera pripadnosti i redak 43 se izvodi.
\end{frame}

% ------------------------------------------------------------
\begin{frame}[fragile]{Test koji prolazi zbog druge provjere (III)}
\vspace{-0.2em}
\begin{lstlisting}[language=python]
# tests/test_registrations.py
async def test_registration_under_other_existing_competition_is_404(
    client, db, lifter, competition
):
    other = Competition(
        name="Drugi kup", location="Split", date=date(2026, 6, 20),
        prelim_deadline=datetime(2026, 5, 20, tzinfo=timezone.utc),
        final_deadline=datetime(2026, 6, 5, tzinfo=timezone.utc),
    )
    db.add(other)
    await db.commit()
    reg = await add_registration(db, lifter, competition, "83")
    headers = await auth_header(client, "testclub", "klub123")
    resp = await client.get(url(other, reg.id), headers=headers)
    assert resp.status_code == 404
    assert resp.json()["code"] == "not_found"
\end{lstlisting}
\end{frame}

% ------------------------------------------------------------
\begin{frame}{Granica mjere}
Pokrivenost broji izvedene retke, a ne provjerene tvrdnje. Test koji pozove funkciju i ništa ne tvrdi daje jednaku pokrivenost kao test koji provjeri ishod.

\vspace{0.7em}
Ni redak koji ostane neizveden nije uvijek propust:

\vspace{0.3em}
\begin{itemize}
  \itemsep0.4em
  \item testovi zamjenjuju \texttt{get\_db} iz \codefile{app/core/deps.py} vlastitom sesijom nad testnom bazom
  \item grane koje hvataju \texttt{IntegrityError} odgovaraju na utrku dvaju zahtjeva između provjere i upisa, koju je u testu teško izazvati
\end{itemize}

\vspace{0.6em}
\begin{konvencija}
U ovom projektu svako pravilo i svako odbijanje u sloju \codefile{services} ima vlastiti test. Mjera pokrivenosti nije cilj, nego služi za suprotno pitanje: koji dio koda nijedan test ne dira.
\end{konvencija}
\end{frame}

% ------------------------------------------------------------
\begin{frame}{Vježba: mjera i test pod istim imenom}
Radite u mapi \codefile{api}, uz pokrenutu bazu i aktivno virtualno okruženje.

\vspace{0.8em}
\begin{enumerate}
  \itemsep0.4em
  \item Pokrenite \texttt{pytest -{}-cov} i zabilježite postotak i neizvedene retke za \codefile{registration\_service.py}.
  \item U \codefile{pyproject.toml} zakomentirajte redak \texttt{concurrency} i ponovite prvi korak.
  \item Vratite redak. U \codefile{registration\_service.py} nađite uvjet koji odbija prijavu drugog natjecanja i uklonite samo tu provjeru.
  \item Pokrenite \texttt{pytest tests/test\_registrations.py} i zabilježite imena testova koji su pali.
\end{enumerate}

\vspace{0.8em}
Zašto se postotak u drugom koraku promijenio iako nijedan test nije promijenjen, i po čemu se u četvrtom koraku razlikuju dva testa prijave pod tuđim natjecanjem?
\end{frame}

% ============================================================
\sekcija{Testovi klijentskog dijela}{Pogreške koje backend ne vidi, vrijeme i zona kao ulaz testa, isti podatak na dvjema stranama, lažne zamjene modula i mreže, provjera testa namjernim kvarom.}

% ------------------------------------------------------------
\begin{frame}{Pogreške frontenda}
Backend odbija zahtjev koji krši pravilo i time štiti podatke. Frontend koji pogriješi podatke ne kvari, ali korisniku ne kaže istinu:

\vspace{0.6em}
\begin{itemize}
  \itemsep0.45em
  \item datum natjecanja ispiše se dan ranije korisniku zapadno od Greenwicha
  \item poruka poslužitelja ne dođe do polja na koje se odnosi, pa obrazac pogrešku ne pokaže
  \item pohrana nakon neuspjelog dohvata ostane u stanju dohvata, pa prikaz zauvijek čeka
  \item istodobni zahtjevi s isteklim tokenom pokrenu svaki svoju obnovu tokena
\end{itemize}

\vspace{0.6em}
Nijednu od tih pogrešaka testovi backenda ne vide, jer se dogode u pregledniku. Zato frontend dobiva vlastite testove.
\end{frame}

% ------------------------------------------------------------
\begin{frame}[fragile]{Alat za jedinične testove}
\vspace{-0.2em}
\begin{lstlisting}[language=typescript]
// vitest.config.ts
process.env.TZ = 'America/Los_Angeles'

export default defineConfig({
  resolve: {
    alias: { '@': fileURLToPath(new URL('./src', import.meta.url)) },
  },
  test: {
    include: ['src/**/__tests__/*.test.ts'],
    environment: 'node',
  },
})
\end{lstlisting}

\vspace{0.2em}
Vitest izvodi testove u Nodeu i pokreće se naredbom \texttt{npm test}. Uz \texttt{environment: 'node'} test nema stablo elemenata, jer ga testirane funkcije ne trebaju. Redak s \texttt{TZ} objašnjava sljedeći slajd.
\end{frame}

% ------------------------------------------------------------
\begin{frame}[fragile]{Zona u kojoj se test izvodi (I)}
Kalendar koji termin pohrani kao ponoć po UTC-u u Zagrebu ga prikaže tog dana u jedan sat, a u Kaliforniji prethodnog dana poslijepodne. Isto se dogodi datumu bez vremena koji se pretvori u \texttt{Date}:

\begin{lstlisting}[language=typescript, keepspaces=true]
new Date('2027-03-14').toLocaleDateString('hr-HR')
// u zoni Europe/Zagreb:      14. 03. 2027.
// u zoni America/Los_Angeles: 13. 03. 2027.
\end{lstlisting}

\vspace{0.4em}
\texttt{new Date('2027-03-14')} je ponoć po UTC-u, a ispis je u zoni računala.

\vspace{0.4em}
Test pokrenut na računalu u Zagrebu prolazi, a organizator na putu zapadno od Greenwicha vidi dan ranije. Pogrešku sakriva zona u kojoj se test izvodi.
\end{frame}

% ------------------------------------------------------------
\begin{frame}[fragile]{Zona u kojoj se test izvodi (II)}
Vitest zato izvodi testove u zoni zapadno od UTC-a, postavkom \texttt{process.env.TZ}. Kod koji miješa zone tada pada na svakom računalu:

\begin{lstlisting}[language=typescript]
// src/lib/format.ts
export function formatDate(value: string): string {
  const match = /^(\d{4})-(\d{2})-(\d{2})/.exec(value)
  if (!match) return value
  const [, year, month, day] = match
  return `${Number(day)}. ${Number(month)}. ${year}.`
}
\end{lstlisting}

\vspace{-0.2em}
\begin{lstlisting}[language=typescript]
// src/lib/__tests__/format.test.ts
it('daje isti dan i zapadno od UTC-a', () => {
  expect(formatDate('2027-03-14')).toBe('14. 3. 2027.')
})
\end{lstlisting}

\vspace{0.2em}
\texttt{formatDate} datum čita kao tekst, pa daje isti dan u svakoj zoni.
\end{frame}

% ------------------------------------------------------------
\begin{frame}[fragile]{Prelazak na ljetno vrijeme}
Rok upisan po zagrebačkom vremenu dobiva pomak \texttt{+01:00} zimi i \texttt{+02:00} ljeti. Zagreb prelazi na ljetno vrijeme 28.03.2027., a Kalifornija dva tjedna ranije:

\begin{lstlisting}[language=typescript]
// src/lib/__tests__/datetime.test.ts
it('<@prati prelazak na ljetno vrijeme u Zagrebu, a ne u zoni računala@>',
  () => {
    expect(fromLocalInput('2027-03-27T12:00'))
      .toBe('2027-03-27T12:00:00+01:00')
    expect(fromLocalInput('2027-03-28T12:00'))
      .toBe('2027-03-28T12:00:00+02:00')
  })
\end{lstlisting}

\vspace{0.3em}
Dana 27.03. u zoni testa ljetno vrijeme već vrijedi, a u Zagrebu još ne. Test zato razlikuje funkciju koja pomak računa za Zagreb od one koja ga računa za zonu računala ili ga uzima kao stalan.
\end{frame}

% ------------------------------------------------------------
\begin{frame}[fragile]{Današnji dan kao argument}
Funkcija koja sama čita sat ne može se provjeriti na granici bez zamrzavanja vremena. Funkcija \texttt{validateLifter}, koja provjerava obrazac natjecatelja, prima današnji dan kao argument, pa test zadaje i dan i granicu:

\begin{lstlisting}[language=typescript]
// src/lib/__tests__/validation.test.ts
const TODAY = '2026-11-26'

it('<@datum rođenja mora biti prije današnjeg dana@>', () => {
  expect(validateLifter(
    { ...VALID, birth_date: '2026-11-25' }, TODAY)).toEqual({})
  expect(validateLifter({ ...VALID, birth_date: TODAY }, TODAY))
    .toEqual({ birth_date: '<@Datum rođenja mora biti u prošlosti.@>' })
})
\end{lstlisting}

\vspace{0.3em}
Test ima oba slučaja uz granicu. Prvi pokazuje da pravilo propušta ispravan unos, drugi da odbija današnji dan.
\end{frame}

% ------------------------------------------------------------
\begin{frame}[fragile]{Isti popis na dvjema stranama (I)}
Provjera s početka predavanja u projektu postoji kao test. Čita datoteku backenda kao tekst i iz nje vadi popis:

\begin{lstlisting}[language=typescript]
// src/lib/__tests__/categories.test.ts
const SCHEMA = fileURLToPath(new URL(
  '../../../../api/app/schemas/registration.py', import.meta.url))

function pythonList(source: string, name: string): string[] {
  const match = new RegExp(`^${name} = \\[(.*)\\]$`, 'm').exec(source)
  if (!match) throw new Error(...)  <@\komts{popis nije pronađen}@>
  return [...match[1]!.matchAll(/"([^"]+)"/g)].map((m) => m[1]!)
}
\end{lstlisting}

\vspace{0.3em}
Regularni izraz traži redak koji počinje imenom popisa, a iz njega uzima vrijednosti u navodnicima.
\end{frame}


% ------------------------------------------------------------
\begin{frame}[fragile]{Isti popis na dvjema stranama (II)}
\begin{lstlisting}[language=typescript]
// src/lib/__tests__/categories.test.ts
const source = readFileSync(SCHEMA, 'utf8')

it('<@ženske kategorije jednake su onima u shemi backenda@>', () => {
  expect([...CATEGORIES.F])
    .toEqual(pythonList(source, 'FEMALE_CATEGORIES'))
})
\end{lstlisting}

\vspace{0.4em}
Uz izmjenu s početka predavanja, kategoriju 43 samo na backendu, ovaj test pada s imenom pravila koje je prekršeno. Komentar \enquote{mijenjati zajedno} sada provjerava stroj.

\vspace{0.4em}
Cijena je ovisnost o obliku retka u Pythonu: popis prelomljen u više redaka izraz ne pronalazi, pa test pada i bez razilaženja.
\end{frame}

% ------------------------------------------------------------
\begin{frame}[fragile]{Lažna zamjena modula (I)}
Pohrana natjecatelja dohvaća popis funkcijom \texttt{fetchLifters}, koja šalje zahtjev. Test pohrane ne treba poslužitelj: modul se zamjenjuje lažnom inačicom (engl. \emph{mock}), čije odgovore zadaje test.

\begin{lstlisting}[language=typescript]
// src/stores/__tests__/lifters.test.ts
vi.mock('@/api/lifters', () => ({ fetchLifters: vi.fn() }))

beforeEach(() => {
  setActivePinia(createPinia())
  vi.mocked(fetchLifters).mockReset()
  useAuthStore().user = { id: 5, username: 'behemot@plreg.hr',
    role: 'club', club_id: 7, is_active: true }
})
\end{lstlisting}

\vspace{0.3em}
\texttt{vi.mock} zamjenjuje modul za svaki uvoz u ovoj datoteci, pa i za uvoz unutar pohrane. \texttt{setActivePinia(createPinia())} prije svakog testa daje novu, praznu pohranu, pa stanje jednog testa ne prelazi u drugi.
\end{frame}

% ------------------------------------------------------------
\begin{frame}[fragile]{Lažna zamjena modula (II)}
\begin{lstlisting}[language=typescript]
// src/stores/__tests__/lifters.test.ts
it('<@neuspjeli dohvat ostavlja prikazani popis i bilježi neuspjeh@>',
  async () => {
    vi.mocked(fetchLifters)
      .mockResolvedValueOnce([ANA])
      .mockRejectedValueOnce(new Error('<@mreža@>'))
    const store = useLiftersStore()
    await store.load()
    await expect(store.load()).rejects.toThrow('<@mreža@>')
    expect(store.items).toEqual([ANA])
    expect(store.failed).toBe(true)
    expect(store.pending).toBe(false)
  })
\end{lstlisting}

\vspace{0.3em}
Prvi dohvat uspije, drugi padne. Popis iz prvoga mora ostati, a neuspjeh biti zabilježen. Zadnja tvrdnja čuva ono što se lako zaboravi: prikaz koji nakon neuspjeha zauvijek čeka.
\end{frame}

% ------------------------------------------------------------
\begin{frame}[fragile]{Lažni poslužitelj}
Kad pristupni token istekne, prvi zahtjev koji dobije \texttt{401} traži novi par tokena, a ostali čekaju ishod iste obnove. Za provjeru toga test treba poslužitelj koji odgovara po pravilu, pa zamjenjuje globalne objekte preglednika:

\begin{lstlisting}[language=typescript]
// src/api/__tests__/http.test.ts
beforeEach(() => {
  vi.stubGlobal('localStorage', memoryStorage())
  server = fakeServer()
  vi.stubGlobal('fetch', server)
})
\end{lstlisting}

\vspace{0.3em}
\texttt{memoryStorage} drži zapise u memoriji, pa test ne ovisi o pohrani okruženja u kojem se izvodi. \texttt{fakeServer} prihvaća samo token \texttt{novi}: na zahtjev za obnovu vraća par s njim, a na svaki drugi token odgovara \texttt{401}. Bilježi i svaki poziv, pa test može prebrojiti zahtjeve za obnovu.
\end{frame}

% ------------------------------------------------------------
\begin{frame}[fragile]{Istodobni odgovori 401}
\begin{lstlisting}[language=typescript]
// src/api/__tests__/http.test.ts
it('tri istodobna odgovora 401 <@pokreću@> jednu obnovu', async () => {
  saveTokens({ access_token: 'stari', refresh_token: 'obnova',
    token_type: 'bearer' })
  await Promise.all([
    request('GET', '/competitions/'),
    request('GET', '/auth/me'),
    request('GET', '/clubs/'),
  ])
  expect(refreshCalls()).toHaveLength(1)
})
\end{lstlisting}

\vspace{0.4em}
Tri zahtjeva kreću s istim tokenom koji poslužitelj ne prihvaća i dobivaju \texttt{401} u istom trenutku. Tvrdnja ne gleda odgovore, nego broj zahtjeva za obnovu koje je lažni poslužitelj primio.
\end{frame}

% ------------------------------------------------------------
\begin{frame}[fragile]{Provjera testa namjernim kvarom (I)}
Test koji od prvog pokretanja prolazi nikad nije pokazao da može pasti: ako tvrdnja ne gleda redak koji test čuva, prolazit će i bez njega.

\begin{lstlisting}[language=typescript]
// src/api/http.ts
let refreshing: Promise<boolean> | null = null
function refreshTokens(): Promise<boolean> {
  refreshing ??= (async () => {
    <@\komts{zahtjev za obnovu i spremanje novog para}@>
    ...
  })().finally(() => {
    refreshing = null
  })
  return refreshing
}
\end{lstlisting}

\vspace{0.3em}
Jednu obnovu čuva operator \texttt{??=}: nova se pokreće samo ako nijedna nije u tijeku.
\end{frame}

% ------------------------------------------------------------
\begin{frame}[fragile]{Provjera testa namjernim kvarom (II)}
Operator \texttt{??=} zamijeni se s \texttt{=}, pa svaki poziv pokreće novu obnovu:

\begin{lstlisting}[keepspaces=true]
 > src/api/__tests__/http.test.ts (6 tests | 1 failed)
     x tri istodobna odgovora 401 <@pokreću@> jednu obnovu
 Test Files  1 failed (1)
      Tests  1 failed | 5 passed (6)
\end{lstlisting}

\vspace{0.4em}
Pao je test koji tu pogrešku čuva, i nijedan drugi. Kvar koji ne obori nijedan test pokazuje da ponašanje koje bi kvar promijenio nitko ne provjerava.

\vspace{0.4em}
Provjera traje jedno pokretanje i radi se rukom, za pravila koja se lako pokvare. Alati za testiranje mutacijama (engl. \emph{mutation testing}) unose takve kvarove sami, redom kroz cijeli kod.
\end{frame}

% ------------------------------------------------------------
\begin{frame}{Opseg jediničnih testova}
Test se piše, čita i mijenja uz svaku izmjenu koda koji dira. Zato se ne piše za:

\vspace{0.6em}
\begin{itemize}
  \itemsep0.45em
  \item ponašanje \emph{frameworka}, primjerice da se prikaz osvježi kad se vrijednost promijeni
  \item tuđe knjižnice, koje imaju vlastite testove
  \item predložak bez ijedne odluke, koji bi test samo prepisao
\end{itemize}

\vspace{0.6em}
Ostaje ono što sadrži odluku: pretvorba podataka, uvjet koji radnju skriva ili onemogućuje i ponašanje nakon neuspjelog zahtjeva.

\vspace{0.6em}
\begin{konvencija}
U ovom projektu jedinične testove imaju čiste funkcije iz \codefile{src/lib}, pohrane i sloj koji šalje zahtjeve. Komponente provjerava test kroz preglednik.
\end{konvencija}
\end{frame}

% ------------------------------------------------------------
\begin{frame}[fragile]{Vježba: zona i namjerni kvar}
Radite u mapi \codefile{web}. Testovi se pokreću naredbom \texttt{npm test}.

\vspace{0.4em}
\begin{enumerate}
  \itemsep0.3em
  \item U \codefile{src/lib/format.ts} zamijenite tijelo funkcije \texttt{formatDate} ovim retcima i pokrenite testove:
\begin{lstlisting}[language=typescript]
const d = new Date(value)
return `${d.getDate()}. ${d.getMonth() + 1}. ${d.getFullYear()}.`
\end{lstlisting}
  \item U \codefile{vitest.config.ts} zakomentirajte redak koji postavlja \texttt{TZ} i ponovite testove.
  \item Vratite oba retka. U \codefile{src/api/http.ts} promijenite redak koji pamti obnovu u tijeku tako da svaki poziv \texttt{refreshTokens} pokrene novu, i ponovite testove.
\end{enumerate}

\vspace{0.4em}
Po čemu se razlikuju popisi palih testova u prvom i drugom koraku i što ta razlika govori o korisniku zapadno od Greenwicha?
\end{frame}

% ============================================================
\sekcija{Test kroz preglednik}{Pogreške u spoju dijelova, polazno stanje i priprema podataka kroz API, lokatori po ulozi i oznaci, provjera pristupačnosti i nalazi na uskom zaslonu.}

% ------------------------------------------------------------
\begin{frame}{Pogreške u spoju dijelova}
Jedinični testovi provjeravaju dijelove. Svaki od njih može prolaziti, a aplikacija svejedno ne radi:

\vspace{0.6em}
\begin{itemize}
  \itemsep0.45em
  \item pohrana prijave sprema token pod jednim ključem, a sloj koji šalje zahtjeve čita ga pod drugim
  \item provjera prije navigacije pita je li korisnik prijavljen prije nego je dohvat korisnika završio
  \item zahtjev bez kose crte na kraju adrese dobije preusmjerenje na adresu koja iza posrednika vodi na krivo mjesto
\end{itemize}

\vspace{0.6em}
Takve pogreške postoje samo u spoju dijelova. Vidi ih test koji aplikaciju pokreće onako kako je pokreće korisnik: u pregledniku, nad pokrenutim backendom i bazom.
\end{frame}

% ------------------------------------------------------------
\begin{frame}[fragile]{Alat za testove kroz preglednik (I)}
\begin{lstlisting}[language=typescript]
// playwright.config.ts
const BASE_URL = process.env.E2E_URL ?? 'http://127.0.0.1:5173'
export default defineConfig({
  testDir: 'e2e',
  retries: 0,
  use: { baseURL: BASE_URL, locale: 'hr-HR',
         timezoneId: 'Europe/Zagreb', trace: 'retain-on-failure' },
  projects: [
    { name: '<@računalo@>', use: { ...devices['Desktop Chrome'] } },
    { name: 'mobitel', use: { ...devices['Pixel 7'] } },
  ],
  ...  <@\komts{pokretanje aplikacije: sljedeći slajd}@>
})
\end{lstlisting}

\vspace{0.3em}
Svaki test izvodi se u pravom pregledniku dvaput, kao na računalu i kao na mobitelu. Zapis tijeka (engl. \emph{trace}) ostaje samo za pali test.
\end{frame}


% ------------------------------------------------------------
\begin{frame}[fragile]{Alat za testove kroz preglednik (II)}
\begin{lstlisting}[language=typescript]
// playwright.config.ts
  webServer: process.env.E2E_URL ? undefined : {
    command: 'npm run dev', url: BASE_URL, reuseExistingServer: true,
  },
\end{lstlisting}

\vspace{0.3em}
Bez varijable \texttt{E2E\_URL} Playwright sam pokreće razvojni poslužitelj frontenda, ako već ne radi, a on module prevodi tek pri prvom zahtjevu. Uz varijablu testovi idu na zadanu adresu, primjerice na izgrađenu aplikaciju.

\vspace{0.5em}
\begin{konvencija}
Projekt test koji ponekad padne popravlja, a ne ponavlja (\texttt{retries: 0}), jer bi ponavljanje nestabilnost sakrilo. Mjerodavno je pokretanje nad izgrađenom aplikacijom.
\end{konvencija}
\end{frame}

% ------------------------------------------------------------
\begin{frame}[fragile]{Test prijave}
\begin{lstlisting}[language=typescript]
// e2e/prijava.spec.ts
test('kriva lozinka daje poruku i ostavlja korisnika na obrascu',
  async ({ page }) => {
    await page.goto('/prijava')
    await expectAccessible(page)  <@\komts{pristupačnost, na kraju cjeline}@>
    await page.getByLabel('<@Korisničko ime@>').fill(CLUB.username)
    await page.getByLabel('Lozinka').fill('kriva-lozinka')
    await page.getByRole('button', { name: 'Prijavi se' }).click()
    await expect(page.getByRole('alert'))
      .toHaveText('<@Neispravno korisničko ime ili lozinka@>')
    await expect(page).toHaveURL(/\/prijava$/)
  })
\end{lstlisting}

\vspace{0.3em}
Test radi ono što radi korisnik: otvara stranicu, upisuje u polja, pritišće gumb i čita poruku. Zadnja tvrdnja provjerava adresu, jer kriva lozinka ne smije voditi dalje.
\end{frame}

% ------------------------------------------------------------
\begin{frame}[fragile]{Lokatori po ulozi i oznaci}
\texttt{getByLabel} nalazi polje po tekstu povezane oznake, a \texttt{getByRole} element po ulozi i nazivu, kako ih vidi čitač zaslona:

\begin{lstlisting}[language=html, morecomment={[s]{<!--}{-->}}]
<!-- src/views/LoginView.vue -->
<label for="username"><@Korisničko ime@></label>
<input id="username" v-model="username"
       autocomplete="username" required />
\end{lstlisting}

\vspace{0.3em}
Test ne nalazi polje čija oznaka s njim nije povezana, jer se \texttt{for} i \texttt{id} razlikuju. Čitač zaslona takvom polju ne čita oznaku, pa je lokator i provjera pristupačnosti.

\vspace{0.4em}
\begin{konvencija}
Elementi se mogu tražiti i po atributu \texttt{data-testid}, koji ne ovisi o tekstu. Projekt ih traži po ulozi i oznaci, pa promjena teksta koji korisnik vidi smije oboriti test.
\end{konvencija}
\end{frame}

% ------------------------------------------------------------
\begin{frame}{Polazno stanje (I)}
Laboratorijski pokus počinje od pripremljenog uzorka, a ne od onoga što je ostalo iza prethodne skupine. Test kroz preglednik radi nad stvarnom bazom, pa ishod ovisi o tome što u njoj zatekne.

\vspace{0.9em}
Test koji odjavi prijavu ostavlja je odjavljenom. Sljedeće pokretanje istog testa uz istu prijavu nalazi gumb \enquote{Vrati prijavu} umjesto gumba \enquote{Odjavi} i pada, iako se kod nije promijenio.

\vspace{0.9em}
Isto se dogodi dvama testovima koji se izvode istodobno i dijele zapis: ishod ovisi o tome koji je prvi stigao.
\end{frame}


% ------------------------------------------------------------
\begin{frame}{Polazno stanje (II)}
\vspace{0.3em}
\begin{center}
{\small
\setlength{\tabcolsep}{6pt}
\renewcommand{\arraystretch}{1.3}
\begin{tabular}{@{}L{4.2cm}L{3.6cm}L{4.0cm}@{}}
  \toprule
  \textbf{Pristup} & \textbf{Postupak} & \textbf{Cijena} \\
  \midrule
  prazna baza za svaki test & brisanje i ponovno punjenje & sporo; test mora dirati bazu mimo API-ja \\
  zajednički zapisi & svi testovi rade nad istim podatcima & test ovisi o tome što su drugi promijenili \\
  vlastiti zapisi & test ih stvara kroz API, s jedinstvenim imenom & baza raste sa svakim pokretanjem \\
  \bottomrule
\end{tabular}
}
\end{center}

\vspace{0.6em}
\begin{konvencija}
U ovom projektu računi dolaze iz početnih podataka, a natjecanje i natjecatelja svaki test stvara sam, s vremenom i slučajnim brojem u imenu.
\end{konvencija}
\end{frame}

% ------------------------------------------------------------
\begin{frame}[fragile]{Priprema podataka kroz API}
\begin{lstlisting}[language=typescript]
// e2e/helpers.ts
export async function createOpenCompetition(
  request: APIRequestContext) {
  const admin = await apiAs(request, ADMIN)
  const competition = await admin('POST', '/competitions/', {
    name: `E2E kup ${unique()}`,
    location: 'Zadar',
    date: at(60).slice(0, 10),
    prelim_deadline: at(20),
    final_deadline: at(40),
  })
  <@\komts{vraća natjecanje i funkcije za pomicanje rokova}@>
  ...
}
\end{lstlisting}

\vspace{0.3em}
Natjecanje nastaje zahtjevom organizatora, mimo obrasca, jer obrazac nije predmet ovog testa.
\end{frame}

% ------------------------------------------------------------
\begin{frame}[fragile]{Rok u testu kroz preglednik (I)}
Test kroz preglednik ne može zamrznuti sat backenda: backend radi u svojem procesu i čita vlastiti sat. Test zato pomiče podatke i rokove natjecanja premješta u prošlost:

\begin{lstlisting}[language=typescript]
// e2e/helpers.ts
const DAY = 24 * 60 * 60 * 1000
const at = (days: number) =>
  new Date(Date.now() + days * DAY).toISOString()

  <@\komts{u povratnoj vrijednosti funkcije createOpenCompetition}@>
  passPrelim: () => admin('PATCH', `/competitions/${competition.id}`,
    { prelim_deadline: at(-2) }),
  close: () => admin('PATCH', `/competitions/${competition.id}`,
    { prelim_deadline: at(-2), final_deadline: at(-1) }),
\end{lstlisting}

\vspace{0.3em}
\texttt{at(n)} je trenutak $n$ dana od sada. Rok se pomiče za dane, a ne za sekunde, pa ishod ne ovisi o tome koliko test traje.
\end{frame}


% ------------------------------------------------------------
\begin{frame}[fragile]{Rok u testu kroz preglednik (II)}
\begin{lstlisting}[language=typescript]
// e2e/prijave.spec.ts
test(
  '<@nakon preliminarnog roka nema nove prijave, a odjava je dopuštena@>',
  async ({ page, request }) => {
    const { competition, passPrelim } =
      await createOpenCompetition(request)
    const lifter = await createLifter(request, competition.id)
    await passPrelim()
    await login(page, CLUB)
    await page.goto(`/natjecanja/${competition.id}`)
    await expect(page.getByText('Rok za nove prijave je istekao.'))
      .toBeVisible()
    ...  <@\komts{gumba za novu prijavu nema, a gumb Odjavi je omogućen}@>
  })
\end{lstlisting}

\vspace{0.3em}
Natjecatelj se prijavljuje dok je razdoblje \texttt{open}, a rok se pomiče tek iza toga. Test tako provjerava odjavu prijave koja je nastala prije roka.
\end{frame}

% ------------------------------------------------------------
\begin{frame}[fragile]{Provjera pristupačnosti}
Smjernice za pristupačnost web sadržaja (engl. \emph{Web Content Accessibility Guidelines}, WCAG) propisuju kriterije za korisnike čitača zaslona, tipkovnice i povećanog prikaza. Dio se može provjeriti strojno: kontrast teksta, oznaka polja, naziv gumba, uloga elementa.

\begin{lstlisting}[language=typescript]
// e2e/helpers.ts
export async function expectAccessible(page: Page) {
  const results = await new AxeBuilder({ page })
    .withTags(['wcag2a', 'wcag2aa', 'wcag21a', 'wcag21aa', 'wcag22aa'])
    .analyze()
  expect(results.violations.map((v) =>
    `${v.id}: ${v.help} (${v.nodes.length})`)).toEqual([])
}
\end{lstlisting}

\vspace{0.2em}
Knjižnica axe provjerava kriterije razina A i AA nad stranicom kakva je u trenutku poziva. Ne može procijeniti je li oznaka smislena ni je li redoslijed fokusa logičan; to ostaje ručnoj provjeri.
\end{frame}

% ------------------------------------------------------------
\begin{frame}[fragile]{Okvir tablice na mobitelu (I)}
U prvom pokretanju testova kroz preglednik test pravila nakon konačnog roka u projektu \texttt{mobitel} pada na provjeri pristupačnosti:

\begin{lstlisting}[keepspaces=true]
[mobitel] > prijave.spec.ts > nakon <@konačnog@> roka radnje su ...
    + Array [
    +   "scrollable-region-focusable: Scrollable region must
         have keyboard access (1)",
    + ]
\end{lstlisting}

\vspace{0.4em}
Isti test u projektu \texttt{računalo} prolazi. Provjera tipova i jedinični testovi prolaze.

\vspace{0.4em}
Na računalu tablica stane u širinu. Na mobitelu se okvir tablice pomiče vodoravno, a korisnik tipkovnice do stupaca izvan zaslona ne dolazi, jer okvir ne može dobiti fokus.
\end{frame}

% ------------------------------------------------------------
\begin{frame}[fragile]{Okvir tablice na mobitelu (II)}
Okvir tablice postaje zasebna komponenta, koju koriste sve tablice:

\begin{lstlisting}[language=html, morecomment={[s]{<!--}{-->}}]
<!-- src/components/TableScroll.vue -->
<script setup lang="ts">
defineProps<{ label: string }>()
</script>
<template>
  <div class="table-wrap" tabindex="0" role="region"
       :aria-label="label">
    <slot />
  </div>
</template>
\end{lstlisting}

\vspace{0.3em}
\texttt{tabindex="0"} okviru daje fokus, pa ga strelice pomiču. \texttt{role="region"} i \texttt{aria-label} čitaču zaslona kažu što je u okviru.

\vspace{0.3em}
Cijena je dodatno zaustavljanje pri kretanju tipkom \kbd{Tab}, i onda kad se okvir ne pomiče.
\end{frame}

% ------------------------------------------------------------
\begin{frame}{Širina stranice na mobitelu (I)}
Test prijave natjecatelja na mobitelu pada na kliku na gumb Odjavi: klik prima obrazac iznad tablice. Snimka stranice pokazuje razlog:

\vspace{0.4em}
\begin{center}
\begin{tikzpicture}
  \node[anchor=south west, inner sep=0] (img) at (0,0)
    {\includegraphics[width=7.2cm]{slike/sirina-stranice-mobitel.png}};
  \fill[red, opacity=0.10] (0.730*7.2, 0) rectangle (7.2, 3.52);
  \draw[dashed, very thick, color=red!75!black]
    (0.730*7.2, -0.15) -- (0.730*7.2, 3.67);
  \node[font=\scriptsize, color=red!75!black, anchor=north east]
    at (0.730*7.2 - 0.05, -0.15) {rub zaslona};
  \node[font=\scriptsize, color=red!75!black, anchor=north west]
    at (0.730*7.2 + 0.05, -0.15) {višak};
\end{tikzpicture}
\end{center}

\vspace{0.3em}
Navigacija završava na rubu zaslona, a stranica se nastavlja desno od njega, iako je svaki vidljivi element uži od zaslona.
\end{frame}

% ------------------------------------------------------------
\begin{frame}[fragile]{Širina stranice na mobitelu (II)}
Gumb u retku tablice nosi skriveni natpis za čitač zaslona, s razredom \texttt{.sr-only}, koji je postavljen apsolutno. Apsolutno postavljen element smješta se prema najbližem pozicioniranom pretku. Okvir tablice to nije bio, pa se natpis smjestio prema cijeloj stranici, izašao iz okvira koji se pomiče i proširio stranicu.

\begin{lstlisting}
/* src/assets/main.css */
.sr-only { position: absolute; ... }
.table-wrap {
  position: relative;
  overflow-x: auto;
  ...
}
\end{lstlisting}

\vspace{0.3em}
Uz \texttt{position: relative} okvir tablice postaje taj predak, pa natpis ostaje u okviru koji se pomiče.
\end{frame}

% ------------------------------------------------------------
\begin{frame}[fragile]{Širina stranice na mobitelu (III)}
Ispravak čuva tvrdnja u testu prijave natjecatelja:

\begin{lstlisting}[language=typescript]
// e2e/helpers.ts
export async function expectNoHorizontalScroll(page: Page) {
  const overflow = await page.evaluate(() =>
    document.documentElement.scrollWidth
      - document.documentElement.clientWidth)
  expect(overflow, '<@stranica je šira od zaslona@>')
    .toBeLessThanOrEqual(0)
}
\end{lstlisting}

\vspace{0.4em}
\texttt{scrollWidth} je širina sadržaja, a \texttt{clientWidth} širina vidljivog dijela. Razlika veća od nule znači da se stranica pomiče vodoravno.

\vspace{0.4em}
Oba kvara prošla su provjeru tipova, jedinične testove i test kroz preglednik na računalu. Uhvatio ih je isti test na uskom zaslonu, jer natjecatelju daje dugo prezime; uz kratka imena natpis ostaje unutar zaslona.
\end{frame}

% ------------------------------------------------------------
\begin{frame}{Pokretanje testova}
\vspace{0.3em}
\begin{center}
{\small
\setlength{\tabcolsep}{6pt}
\renewcommand{\arraystretch}{1.3}
\begin{tabular}{@{}L{3.3cm}L{3.5cm}L{4.8cm}@{}}
  \toprule
  \textbf{Razina} & \textbf{Trajanje} & \textbf{Pokreće se} \\
  \midrule
  jedinični, frontend & sekunde & pri svakoj izmjeni \\
  nad API-jem & desetci sekundi & prije svakog \texttt{commita} \\
  kroz preglednik & desetci sekundi i više & prije spajanja u glavni \texttt{branch} \\
  \bottomrule
\end{tabular}
}
\end{center}

\vspace{0.6em}
Test koji se ne izvodi sam ne izvodi se dugo. Pokretanje se zato veže uz radnju koja se ionako radi, a ne uz nečiju odluku da je vrijeme za provjeru.

\vspace{0.6em}
Kad se svi testovi pokreću sami na svaku izmjenu u repozitoriju, nijedna razina ne ovisi o tome je li se netko sjetio pokrenuti je.
\end{frame}

% ------------------------------------------------------------
\begin{frame}{Vježba: nalazi na uskom zaslonu}
Backend radi nad bazom s početnim podatcima. Testovi se pokreću iz mape \codefile{web}.

\vspace{0.7em}
\begin{enumerate}
  \itemsep0.35em
  \item Pokrenite \texttt{npx playwright test} i zabilježite u kojim se projektima testovi izvode.
  \item U \codefile{src/components/TableScroll.vue} uklonite atribut koji okviru daje fokus i pokrenite \texttt{npx playwright test -{}-project=mobitel}.
  \item Pokrenite iste testove s \texttt{-{}-project=računalo}.
  \item Vratite atribut. U \codefile{src/assets/main.css} nađite deklaraciju koja apsolutno postavljene elemente drži u okviru tablice, uklonite je i ponovite drugi korak.
\end{enumerate}

\vspace{0.7em}
Po čemu se razlikuju ishodi drugog i trećeg koraka i koji element u četvrtom koraku određuje širinu stranice?
\end{frame}

% ------------------------------------------------------------
\begin{frame}{Sažetak}
\begin{itemize}
  \itemsep0.2em
  \item Prolazni test pokazuje ishod za zadane ulaze, ne ispravnost koda.
  \item Niža razina testa pada brže i točnije, a viša vidi spoj dijelova.
  \item Pokrivenost broji izvedene retke, a ne provjerene tvrdnje.
  \item Mjera se provjerava retkom za koji se zna da ga test izvodi.
  \item Neizveden redak otkriva test koji prolazi zbog druge provjere.
  \item Jedinični test frontenda izvodi se u zoni zapadno od UTC-a.
  \item Funkcija koja ovisi o današnjem danu prima ga kao argument.
  \item Podatak na objema stranama čuva test koji čita oba izvora.
  \item Test se provjerava namjernim kvarom koji ga mora oboriti.
  \item Test kroz preglednik stvara svoje podatke i traži po ulozi i oznaci.
  \item Test na uskom zaslonu nalazi kvarove koje računalo ne pokazuje.
\end{itemize}
\end{frame}

\end{document}
